\chapter{Assessment Interview Guide}\label{interview_guide}

This appendix contains the interview guide used for the assessment interviews, which were conducted to classify misconceptions about \gls{llms} among the participants.

\section*{Spørsmål til Forhåndsintervju, LLM-explainer}

\textbf{Forklar med egne ord hva AI er og hvordan det fungerer}
\begin{itemize}
\item Etter dette spørmålet skal AI i denne sammenhengen være forstått som språkmodeller, og ikke AI generelt.
\item Dersom det er naturlig, still konkrete oppfølgingsspørmål for å dekke de ulike kategoriene:
	\begin{itemize}
	\item Hvordan oppererer språkmodeller? (OVERAUTO, REALTIME, CONTEXNEG, SEMUNDER, SOURCEATTR, CONDTREAT)
	\item Til hvilken grad stoler du på en språkmodells forståelse og svar? (OVERACC, HALLBLIND, SEMUNDER)
	\item Hvor kommer informasjonen fra når du bruker en språkmodell, er det noen problemstillinger rundt dette og hvem som eier denne informasjonen? (SOURCEATTR, PLAGBLIND, COPYBLIND)
	\end{itemize}
\end{itemize}

\textbf{Forklar hvordan du bruker språkmodeller i hverdagen, og hva du bruker dem til}
\begin{itemize}
\item Bruk bruksområdene som blir nevnt til å spørre om hvordan denne prosessen fungerer, og hva de tror skjer ``bak kulissene'' når de bruker språkmodeller.
\item Spør om en spesifik prompt de kunne brukt
\item Prøv også å finne ut av til hvilken grad de stoler på språkmodeller, og hva de gjør for å verifisere informasjonen de får fra dem. Også om det er noen situasjoner der de ikke stoler på språkmodeller, og hvorfor.
\end{itemize}

\textbf{Hvordan påvirker språkmodeller deg i hverdagen, både positivt og negativt?}
\begin{itemize}
\item Få vedkommende til å drøfte hvilke effekter dette har
\end{itemize}

\textbf{Hvordan påvirker språkmodeller samfunnet generelt, både nå og i fremtiden?}
\begin{itemize}
\item Få vedkommende til å drøfte hvilke effekter dette har, både på kort og lang sikt, og både positivt og negativt.
\end{itemize}
